% ── §7f Methods: Formal Verification and Reproducibility ──
\section{Formal Verification as a Reproducibility Guarantee}
\label{sec:methods-reproducibility}

\subsection{The Reproducibility Crisis in Computational Biology}
Computational biology faces well-documented challenges with reproducibility:
results may depend on undocumented parameter choices, unstated preprocessing
steps, or subtle implementation details that are lost when code is transferred
between researchers.  In the specific context of contact prediction, a change
as simple as switching from C$\beta$ to C$\alpha$ distance, or from separation
$\geq 6$ to $\geq 12$, can alter reported precision by several percentage
points.

Our approach addresses this by formalising all load-bearing definitions in
Lean~4, where they are checked by the type system and kernel rather than by
human review.  When we state that ``the Gray-code transformation maps
consecutive integers to adjacent codewords,'' this is not an empirical claim
subject to interpretation but a machine-checked theorem whose proof is
verified by the Lean kernel.

\subsection{Scope of Formal Guarantees}
It is important to note what our formal verification does and does not prove.
We prove that the encoding layer correctly implements its specification:
that Gray codes are bijective on their domain, that adjacent integers map to
adjacent codewords, that implicant covers are sound and complete, and that
padding does not corrupt decoded rules.  We do \emph{not} prove that these
specifications are optimal for contact prediction; optimality is an empirical
question that must be settled by experiment.

This distinction mirrors the relationship between compiler correctness proofs
and program performance: a verified compiler guarantees that the output
faithfully implements the source code, but says nothing about whether the
source code is efficient.  Similarly, our verified encoding guarantees that
the pipeline faithfully implements its design, while the design's effectiveness
is established empirically.

\subsection{Verification Infrastructure}
The verification infrastructure consists of two Lean~4 packages.  The first,
\texttt{is\_kmap\_possible}, contains 106 theorems across five modules
covering Gray-code fundamentals (bijectivity, adjacency, involution),
$k$-mer indexing (injectivity and boundedness for $k = 1$ through $4$),
amino-acid encodings (injectivity, within-group distance bounds, charge-group
adjacency), contact-map completeness (symmetry, irreflexivity, Boolean
minimisation semantics), and encoding equivalence (BLOSUM bound, distance
distribution invariance).  The second, \texttt{ContactCircuits}, contains
12 theorems covering the specific assertions made by the circuit-extraction
pipeline: cover completeness, cover soundness, padding safety, separation-bin
monotonicity, and cell-index injectivity.

All proofs use only \texttt{Init} (the Lean standard library) and bounded-domain
tactics (\texttt{decide}, \texttt{native\_decide}).  No theorem depends on
\texttt{sorry} or on axioms beyond the standard logical framework.  The axiom
audit performed via \texttt{\#print axioms} confirms that every theorem
depends at most on \texttt{propext}, \texttt{Quot.sound}, and the
\texttt{native\_decide} evaluation axiom.
